\documentclass[11pt]{article}
\usepackage{amsmath,amssymb,amsthm,hyperref}
\newtheorem{theorem}{Theorem}
\newtheorem{lemma}{Lemma}
\title{A sharp logarithmic barrier for approximate exponential moments}
\author{Research note; independently reviewed}
\date{30 September 2026}
\begin{document}\maketitle
The endpoint lower bound $K\ge c m\log m$ uses a positive measure whose
normalized moments approximate those of $e^{-x}\,dx$. This note shows
that moment accuracy and positivity alone cannot yield an endpoint
atom bound of a smaller order than the reciprocal logarithm.

\begin{lemma}\label{lagbound}
For the ordinary Laguerre polynomials
\[
 L_j(y)=\sum_{k=0}^j(-1)^k\binom jk\frac{y^k}{k!},
\]
one has $|e^{-y/2}L_j(y)|\le1$ for $y\ge0$.
\end{lemma}
\begin{proof}
Here is a self-contained Gaussian proof of this standard inequality.
Give entire functions the norm
\[
 \|f\|^2=\frac1\pi\int_{\mathbb C}|f(z)|^2e^{-|z|^2}\,dA(z).
\]
The monomials $f_j(z)=z^j/\sqrt{j!}$ are orthonormal, by polar
integration. For a complex number $a$, the map
\[
 (D_af)(z)=e^{-|a|^2/2+az}f(z-\bar a)
\]
is an isometry: substituting $u=z-\bar a$ in its squared norm cancels
all exponential factors except $e^{-|u|^2}$.
The coefficient of $z^j$ in $e^{az}(z-\bar a)^j$ is
$\sum_{k=0}^j\binom jk(-|a|^2)^k/k!=L_j(|a|^2)$.
Orthogonality of monomials, or direct integration, therefore gives
\[
 \langle f_j,D_af_j\rangle=e^{-|a|^2/2}L_j(|a|^2).
\]
Cauchy--Schwarz bounds its modulus by one. Choose $|a|^2=y$.
\end{proof}

\begin{theorem}\label{barrier}
For every positive even integer $d$ there is a probability measure
$\nu_d$ on $[0,\infty)$ with an atom of mass exactly $1/(3d)$ at zero,
such that for every integer $k\ge0$,
\[
 0\le\int\frac{x^k}{k!}\,d\nu_d(x)-1
 \le\frac{2^{-d}}{3d}.
\]
The difference is zero for $k<d$. In particular, for every sufficiently
small $\delta>0$ there is a positive probability measure with all
normalized-moment errors at most $\delta$ and an atom at zero of mass
at least $c/\log(1/\delta)$, for an absolute $c>0$.
\end{theorem}
\begin{proof}
Set
\[
 h_d(x)=3e^{-3x}\sum_{j=0}^{d-1}L_j(3x),\qquad a=\frac1{3d},
\]
and define
\[
 d\nu_d(x)=\bigl(e^{-x}-a h_d(x)\bigr)\,dx+a\delta_0(dx).
\]
By Lemma~\ref{lagbound},
\[
 |h_d(x)|\le3d e^{-3x/2},\qquad
 e^{-x}-a h_d(x)\ge e^{-x}-e^{-3x/2}\ge0.
\]
Thus the measure is positive, independently of the parity of $d$.
The continuous part has no atom, so its atom at zero has mass exactly
$a$.

For $z<3$, termwise integration of each finite polynomial gives
\[
 \int_0^\infty e^{zx}\,3e^{-3x}L_j(3x)\,dx
 =\frac{3}{3-z}\left(\frac{-z}{3-z}\right)^j.
\]
Summing the finite geometric series yields
\[
 H_d(z):=\int_0^\infty e^{zx}h_d(x)\,dx
 =1-\left(\frac{-z}{3-z}\right)^d.
\]
In particular $H_d(0)=1$, so $\nu_d$ has total mass one.
For $z<1$ its moment generating function is
\[
 \int e^{zx}\,d\nu_d(x)
 =\frac1{1-z}+a\left(\frac{-z}{3-z}\right)^d.
\]
The correction has sign $(-1)^d$. For even $d$, its Taylor
coefficients are zero below degree $d$, and at degree $k\ge d$ they
are
\[
 a\binom{k-1}{d-1}3^{-k}\ge0.
\]
Their sum is the value at $z=1$ of the correction series, namely
$a2^{-d}$, so each coefficient is at most that quantity. These
coefficients are precisely the normalized-moment errors.
For odd $d$ the same construction is positive but all these errors
have the opposite sign; choosing even $d$ gives the displayed one-sided
version. Finally choose the least even $d\ge\log_2(1/\delta)$.
Then $a2^{-d}\le\delta$, while $a\ge c/\log(1/\delta)$.
\end{proof}

If a finite-support example is desired for any prescribed finite
moment range $0\le k\le N$, replace the positive continuous part by
its positive Gaussian quadrature rule of degree at least $N$, keeping
the atom at zero. Its density is positive for every $x>0$, has all
moments finite, and has infinite support, so the standard orthogonal
polynomial construction applies. The resulting finite positive measure
has exactly the same moments through degree $N$ and the same atom.

Taking $\delta=m^{-1/2}$ shows that even control of every normalized
moment with error $m^{-1/2}$ permits an endpoint atom of order
$1/\log m$. This input is stronger than the available bounds
$(2k+1)/\sqrt m$ for $k\le\sqrt m$. Thus no argument using only those
moment inequalities and positivity can force a power-law improvement
in the endpoint atom bound. Additional information from the original
monotone comparison tensor is necessary. The measures here are not
claimed to arise from such tensors or from fair dice.
\end{document}
